% ═══════════════════════════════════════════════════════════════
% QR-CODE: Lernpfad Reibungskräfte
% Physik Klasse 7 | UE 05
% ═══════════════════════════════════════════════════════════════

\documentclass[a4paper, 12pt]{article}

\usepackage[utf8]{inputenc}
\usepackage[T1]{fontenc}
\usepackage[ngerman]{babel}

\usepackage{helvet}
\renewcommand{\familydefault}{\sfdefault}

\usepackage[margin=2cm]{geometry}
\usepackage{qrcode}
\usepackage[hidelinks]{hyperref}
\usepackage{xcolor}
\usepackage{graphicx}

\definecolor{physik}{HTML}{2c5aa0}

\setlength{\parindent}{0pt}

% URL definieren
\newcommand{\lpurl}{https://devbox42.github.io/sciencesim/physik/kl07/02-kraefte/05-reibungskraefte/LP-05-reibungskraefte.html}

\begin{document}

\pagestyle{empty}

\begin{center}

\vspace*{2cm}

{\LARGE\bfseries\textcolor{physik}{Lernpfad: Reibungskräfte}}

\vspace{0.5cm}

{\large Physik Klasse 7 | Kräfte | UE 05}

\vspace{2cm}

% QR-Code als klickbarer Link
\href{\lpurl}{%
    \qrcode[height=6cm]{\lpurl}%
}

\vspace{1.5cm}

% Vollständiger Link als klickbarer Text
{\large\textbf{Link:}}

\vspace{0.3cm}

\href{\lpurl}{%
    \textcolor{physik}{\texttt{\footnotesize\lpurl}}%
}

\vspace{2cm}

\hrule

\vspace{0.5cm}

{\small Scanne den QR-Code oder klicke auf den Link, um den Lernpfad zu öffnen.}

\end{center}

\end{document}
